\begin{proposition}\label{prop:octic-C75}
There is no point of order $75$ on an elliptic curve over an octic
number field.
\end{proposition}
\begin{proof}
Put $X=X_1(75)$. The complete winding calculation has weighted
dimension $145$, equal to the genus of $X$. Every even-character
component is generated by its winding element. The standard modular
rank-zero argument therefore gives $\operatorname{rk}J_X(\mathbb Q)=0$.
Let $D$ be the rational effective degree-$8$ trace divisor arising
from a point of order $75$.

We use $Y=X/\langle\delta_{26}\rangle$. This quotient parametrizes
an elliptic curve with a point of order $25$ and a cyclic subgroup
of order $3$. The involution fixes the $25$-point, so a fixed
noncuspidal point would require a nonidentity elliptic automorphism
fixing a point of order $25$, which is impossible. Its action on
all $112$ geometric cusps is free as well. Thus $X\to Y$ is etale,
also over $\mathbb Z_2$, and $Y$ has genus $73$ and good reduction
at $2$. Its rational Jacobian has rank zero by pushforward.
The image of a standard rational cusp supplies a rational cusp on $Y$.

For clarity, the reduced model is constructed independently over
$X_1(25)$. In characteristic two any ordinary elliptic curve has
canonical coordinates
\[
 y_c^2+x_cy_c=x_c^3+d x_c^2+1/j.
\]
For a cyclic subgroup of order $3$ its nonzero points have common
$x_c$-coordinate satisfying
\[
 x_c^4+x_c^3+1/j=0.                                   \tag{1}
\]
The certificate verifies the coordinate identity and the division
polynomial exactly on the universal $25$-point family. The resulting
quartic extension is irreducible and has genus $73$, as required.
For efficient computation $X_1(25)$ is first expressed as a degree-$5$
extension using a modular unit; the defining resultant and inverse
coordinate are verified exactly. No bound from an incomplete large
model computation is used.

The Frobenius polynomial gives noncuspidal closed-point counts
\[
 (0,0,0,0,0,0,20,0)
\]
on $X_{\mathbb F_2}$ in degrees one through eight. The cusp signature
is $1^{20},2^{10},4^8,20^2$, checked independently from primitive
Tate cusp labels. Consequently $\overline D$ is entirely cuspidal,
or is a closed degree-$7$ noncuspidal point together with one rational
cusp. All $20$ rational source cusps lift over $\mathbb Q$.

There is one diamond orbit of the $20$ noncuspidal degree-$7$ points.
Indeed the complete ordinary scan over $\mathbb F_{128}$ gives seven
qualifying twists, with $j$-invariants forming one Frobenius orbit.
Each has group $C_{150}$, since its order is $150$ and its smaller
invariant factor divides $127$. Its generators of the order-$75$
subgroup yield precisely the $20$ closed points, on which the diamond
group acts transitively. Their images on $Y$ have degree $7$ and
again form one diamond orbit.

Write $A_7=T_7-1-7\delta_7$. Odd-prime torsion injectivity and
Eichler--Shimura imply $A_7J_Y(\mathbb Q)=0$. The exact boundary
formula gives $A_7c=0$ for the image of any rational source cusp.
Thus the noncuspidal component of a putative $\overline D$ must have
principal $A_7$-image on $Y_{\mathbb F_2}$. For one representative
$P$ the actual isogeny computation gives
\[
 h^0(T_7P-P-7\delta_7P)=0.
\]
Here $T_7P$ has degree $56$ and consists of four degree-$14$ closed
places, each with multiplicity one. The untouched full-division-field
isogeny routine gives exactly twice this divisor when all eight
geometric kernels are listed. Every source point, carried level
structure, proper residue field and norm multiplicity is checked.
An independent finite-field replay verifies all actual moduli records
and the complete degree-$7$ coverage. This excludes the noncuspidal
reduction case.

It remains to consider entirely cuspidal reduction. The closed cusp
signature on $Y_{\mathbb F_2}$, recorded as (degree,width), is
\[
 (1,25)^{10},\quad(1,75)^{10},\quad
 (4,5)^2,\quad(4,15)^2,\quad(10,1),\quad(10,3).
\]
At each width-$25$ rational cusp the source cusp has degree two,
so the multiplicity of $\pi_*\overline D$ is even. All other cusps
of degree at most eight have no further restriction. These facts
are verified by the full primitive-label cusp permutation, including
all residue-degree factors. Hence the number of admissible divisors is
\[
 [t^8](1-t)^{-10}(1-t^2)^{-10}(1-t^4)^{-4}=131790.
\]

Let $\mathcal C$ denote Cartier on
$\Omega=H^0(Y_{\mathbb F_2},\Omega^1)$. The cotangent image of
$T_2-\delta_2$ is
\[
 U=(\mathcal C-1)\Omega,\qquad \dim U=71.
\]
The matrix of Cartier has rank $69$. The diamond $\delta_2$ has order
$10$ on $Y$; on the model (1) it acts on the $25$-point base and fixes
$x_c$. It commutes with Cartier. Its four orbits of relevant closed
cusps have lengths $10,10,2,2$, respectively for widths $25,75,5,15$.
These permutations are independently matched to the complete Tate
cusp action.

We compute sixteen spaces $H^0(\Omega(-mP))$ at four orbit
representatives and transport them by the verified diamond action.
This gives all $128$ required local restriction blocks, including
the nonreduced divisors. Restriction from $U$ has rank eight at
every one of the $131790$ admissible divisors. Independent binary
linear algebra verifies the original kernels, all transports,
projections and exhaustive rank tests.

The proof of Lemma~\ref{lem:quotient-Cartier}, applied at these
admissible pushforwards, now forces $\pi_*D$ to equal its cuspidal
lift. In particular it has cuspidal support, contradicting the
noncuspidal support of $D$. This proves the proposition.
\end{proof}
